This chapter fixes the objects of Section~1 of the paper -- the induction
operator on the flat three-torus and the growth rate $\gamma(u,\eps)$ --
together with the axial Fourier reduction used by the formalisation, and states
the main theorem, Theorem~1.2 of the paper (\texttt{thm:main}), together with
the Lean statement layer that carries it.

Theorem~1.2 of the paper is, up to the deviations recorded after Theorem~\ref{thm:main}, the
Lean predicate \texttt{FastDynamo.MainTheoremComplete}, proved by \texttt{FastDynamo.mainTheoremComplete\_holds}
(Theorem~\ref{thm:main-complete}); it has no hypothesis.  Its \emph{spectral half} is: one
real, divergence-free, axially independent
$\W^{1,\infty}$ velocity field, two constants $\eps_0,\gamma_0>0$ fixed before
the diffusivity, and, for every $0<\eps\leq\eps_0$, a non-zero axial Fourier
index $K_\eps$, an eigenvalue $\lambda_\eps$ with $\re\lambda_\eps\geq\gamma_0$
and a complex magnetic eigenmode
$V^\eps=\e^{2\pi iK_\eps x_3}b^\eps$ of the induction operator, together with
the non-zero real datum $\re V^\eps$ and an explicit real, divergence-free
strong solution obeying the exact exponential norm law.  Its existential form is the Lean
theorem \texttt{FastDynamo.mainPhysicalConclusion\_holds} (alias
\texttt{FastDynamo.main\_theorem\_spectral}).

The route by which that theorem is reached is the paper's, and it is short:
a two-component local seed on the plane (Proposition~\ref{prop:local-seed}), a multiscale
velocity on the torus (Lemma~\ref{lem:global-velocity}), the horizontal operator
$A_{\eps,K}$ on $\Hs^2(\T^2;\CC^2)$ with \emph{no} divergence constraint
(Lemma~\ref{lem:torus-sector-basic}) and its rescaling to the seed
(Lemma~\ref{lem:exact-renormalization}), one weighted-test coercivity estimate for the
background (Proposition~\ref{prop:background-resolvent}), a two-summand parametrix
(Proposition~\ref{prop:global-spectral-transfer}), and the algebraic recovery of the third
magnetic component (Lemma~\ref{lem:axial-recovery}).

The \emph{dynamical} half of Theorem~\ref{thm:main} is the last sentence of the paper's
statement: the topological entropy of every time map of the particle flow vanishes, and the
ideal induction group has exponential growth rate zero in $\Lp^2_\sigma$ operator norm.  It is
the Lean predicate \texttt{FastDynamo.MainTheoremDynamicalClauses}, discharged at the
constructed velocity by \texttt{FastDynamo.mainTheoremDynamical\_holds}; the flow is written
down in closed form and the ideal group is built from the Cauchy formula, in the
particle-dynamics chapter.  The paper proves those two clauses
from Proposition~\ref{prop:particle-flow} and Corollary~\ref{cor:ideal-growth}, which say
more: the linear Lipschitz bound on the flow maps and the two-sided per-datum bounds on the
ideal group.  Lean keeps that surplus in the larger bundle
Proposition~\ref{prop:main-lagrangian}, from which the theorem's two clauses are a projection.
The endpoint that joins the two halves is Theorem~\ref{thm:main-complete}.

The growth-rate clause of the paper's statement, $\gamma(u,\eps)\geq\gamma_0$, is
formalised by the paper's argument (the growing mode is an exponential minorant of the
semigroup norm).  The semigroup on
$\Lp^2_\sigma(\T^3)$ is constructed in Lean from a Hille--Yosida ray, one axial sector at a
time and then by restriction to the solenoidal subspace
(Theorem~\ref{thm:dynamics:solenoidal-semigroup}), and its generator is identified with the
paper's $L_{\eps,u}$ with domain $\Hs^2\cap\Lp^2_\sigma$
(Theorem~\ref{thm:dynamics:solenoidal-generator}), so it is the paper's $\e^{tL_{\eps,u}}$; its
growth bound is \texttt{FastDynamo.gamma}, which is a limit for every $(u,\eps)$
(Lemma~\ref{lem:main:growth-rate-limit}), and the clause is
Theorem~\ref{thm:main-growth-rate}.  The whole of the paper's Lemma~A.1
(\texttt{lem:semigroup-spectral-growth}) --- analyticity, compact resolvent, immediate
compactness, spectral mapping, and $\gamma(u,\eps)$ as the attained spectral bound --- is
Lemma~\ref{lem:semigroup-spectral-growth}, although Theorem~\ref{thm:main} uses none of it.
``Its solution'' is backed by a uniqueness theorem for the curl-form equation
(Lemma~\ref{lem:dynamics:physical-uniqueness}), and $\mathcal T_0$ is \emph{the} solution group
of the ideal equation in weak form (Corollary~\ref{cor:ideal-growth}).

The remarks and definitions of Section~1 of the paper have nodes: the definitions of a
dynamo, a fast dynamo and a slow dynamo (Definition~\ref{def:main:dynamo}, with
Corollary~\ref{cor:main:fast-dynamo}); Remark~1.4, that the real part of the eigenmode is not
an eigenmode (Lemma~\ref{lem:finish:real-part-not-eigenmode}); Remark~1.5
(\texttt{rem:scope-intro}), the zero mean of the magnetic fields
(Lemma~\ref{lem:finish:zero-mean}) and the regularity of the velocity
(Lemma~\ref{lem:main:velocity-regularity}); Remark~1.6, the normalisation
(Lemma~\ref{lem:main:normalisation}); Remark~1.7, the vanishing-diffusivity limit
(Theorem~\ref{thm:vanishing-diffusivity}); and the energy identity
\texttt{eq:energy-competition} (Lemma~\ref{lem:dynamics:energy-identity}).  Conjecture~1.1
(\texttt{conj:smooth-fast-dynamo}) and Remark~1.3 (\texttt{rem:subsequent-resolution}), which
discuss the literature, have no node and no Lean counterpart; the contrast that Remark~1.3
draws for the velocity constructed here is the ideal--resistive gap recorded after
Theorem~\ref{thm:main}.

\section{The induction operator and the statement layer}

\begin{definition}[The induction operator and its axial reduction]
  \label{def:main:induction-operator}
  \lean{FastDynamo.TransverseRealW1Inf, FastDynamo.IsTransverseDivergenceFree,
    FastDynamo.TransverseC3H2, FastDynamo.productPhysicalInductionOperator,
    FastDynamo.productInductionOperator, FastDynamo.reducedInductionOperator,
    FastDynamo.axialLiftH2, FastDynamo.productDivergenceFromH2,
    FastDynamo.DomainOperator.IsEigenpair,
    FastDynamo.DomainOperator.IsStrongSolution}
  \leanok
  For a divergence-free $u\in\W^{1,\infty}(\T^3;\R^3)$ and $\eps>0$, the
  induction operator is
  \[
    L_{\eps,u}B:=\curl(u\times B)+\eps\Delta B
  \]
  on $\Lp^2_\sigma(\T^3)=\{B\in\Lp^2(\T^3;\CC^3):\Div B=0\}$, with domain
  $\Hs^2(\T^3;\CC^3)\cap\Lp^2_\sigma(\T^3)$.  The formalisation also uses the
  auxiliary operator on all vector fields,
  \[
    \mathcal{A}_\eps B
    =\eps\Delta B-(u\cdot\nabla)B+(B\cdot\nabla)u,
    \qquad
    D(\mathcal{A}_\eps)=\Hs^2(\T^3;\CC^3),
  \]
  which agrees with $L_{\eps,u}$ on divergence-free fields.  For a velocity
  independent of $x_3$ every fixed-$K$ axial Fourier subspace is invariant
  under $\mathcal{A}_\eps$: writing $B(y,x_3)=\e^{2\pi iKx_3}b(y)$ and
  $u=(u_\perp,u_3)$,
  \begin{equation}
    A_{\eps,K}b
    =\eps\bigl(\Delta_\perp-(2\pi K)^2\bigr)b
    -u_\perp\cdot\nabla_\perp b-2\pi iKu_3b
    +(b_\perp\cdot\nabla_\perp)u,
    \qquad
    D(A_{\eps,K})=\Hs^2(\T^2;\CC^3).
    \label{eq:main:torus-sector}
  \end{equation}
  The horizontal block of \eqref{eq:main:torus-sector}, which acts on $b_\perp$
  and does not involve $b_3$, is the paper's horizontal operator
  $A_{\eps,K}$ on $\Lp^2(\T^2;\CC^2)$ with domain $\Hs^2(\T^2;\CC^2)$
  (\texttt{eq:torus-sector}).
\end{definition}

In Lean the three operators are
\texttt{FastDynamo.productPhysicalInductionOperator u eps} ($L_{\eps,u}$),
\texttt{FastDynamo.productInductionOperator u eps} ($\mathcal{A}_\eps$) and
\texttt{FastDynamo.reducedInductionOperator u eps K} ($A_{\eps,K}$), and the
lift $b\mapsto\e^{2\pi iKx_3}b$ is
\texttt{FastDynamo.axialLiftH2 K}.  Five points of the encoding matter for
reading every later node.

\begin{itemize}
  \item \emph{The auxiliary operator and the three-component reduction are the
    formalisation's.}  The paper has no auxiliary operator: it works with
    $L_{\eps,u}$ on $\Lp^2_\sigma$ (\texttt{eq:physical-generator}) and, for the
    spectral construction, with the two-component horizontal operator $A_{\eps,K}$ of
    its \texttt{eq:torus-sector}, recovering the axial magnetic component afterwards
    (Lemma~8.1 of the paper, \texttt{lem:axial-recovery}).  Lean's
    \texttt{reducedInductionOperator} is the three-component operator
    \eqref{eq:main:torus-sector}, whose horizontal block is the paper's $A_{\eps,K}$,
    and the passage from the reduced to the physical eigenpair goes through
    $\mathcal{A}_\eps$.
  \item Each operator is a \texttt{FastDynamo.DomainOperator}: a pair of
    bounded maps out of a Banach \emph{graph} domain, an injective inclusion
    into the ambient space and the operator itself.  Thus
    $\W^{1,\infty}(\T^2;\R^3)$, $\Hs^2(\T^2;\CC^3)$ and $\Hs^2(\T^3;\CC^3)$ are
    realised as the closed Fourier-graph submodules
    \texttt{FastDynamo.TransverseRealW1Inf},
    \texttt{FastDynamo.TransverseC3H2} and \texttt{FastDynamo.ProductC3H2},
    complete by construction.  An eigenpair is
    \texttt{DomainOperator.IsEigenpair}, that is, the ambient image of the
    eigenvector is non-zero and the operator acts on it by the scalar; a
    strong solution is \texttt{DomainOperator.IsStrongSolution}, that is, the
    ambient curve is differentiable in time with derivative the operator
    applied to the state.
  \item The velocity type is \emph{transverse}: it is a $\W^{1,\infty}$ field
    on $\T^2$ with values in $\R^3$, so axial independence is part of the type
    rather than an extra hypothesis.  Since the paper's construction also
    produces an $x_3$-independent field, the Lean existential ranges over a
    smaller class than the paper's and is therefore slightly
    \emph{stronger}.  The identification of $\W^{1,\infty}$ with the Lipschitz
    class is not itself formalised; $\W^{1,\infty}$ here means the Fourier
    graph of an essentially bounded value together with two essentially bounded
    weak derivatives, and divergence-freeness is the equation
    \texttt{IsTransverseDivergenceFree u} on the two diagonal weak
    derivatives (there is no axial term because $\partial_3u=0$).
  \item The statement layer has no space $\Lp^2_\sigma$ (the closed solenoidal
    subspace of $\Lp^2(\T^3;\CC^3)$ enters only with the semigroup,
    Definition~\ref{def:main:growth-rate}).  Solenoidality of a field is the
    side condition \texttt{productDivergenceFromH2 V = 0}, and the operators are
    defined on all vector fields, as $\mathcal{A}_\eps$ is.
  \item $L_{\eps,u}$ is realised weakly: \texttt{productPhysicalInductionOperator}
    maps $\Hs^2$ into a Fourier $\Hs^{-2}$ space, because $\curl(u\times B)$ with
    $u\in\W^{1,\infty}$ is a priori only a distribution.  The statement layer
    therefore asserts \emph{both} the $\Lp^2$-valued eigenvalue equation for
    $\mathcal{A}_\eps$ and the $\Hs^{-2}$-valued one for $L_{\eps,u}$; together
    with solenoidality these are equivalent to the paper's $\Lp^2$
    display, which the physical clause alone would not be.
\end{itemize}

\begin{definition}[The maximal sustained growth rate]
  \label{def:main:growth-rate}
  \lean{FastDynamo.gamma, FastDynamo.literalSolenoidalSubspace,
    FastDynamo.literalInductionSemigroup,
    FastDynamo.StronglyContinuousSemigroup.growthBound}
  \leanok
  \uses{def:dynamics:growth-bound, thm:dynamics:solenoidal-semigroup}
  For a time-independent divergence-free $u\in\W^{1,\infty}(\T^3;\R^3)$ and
  $\eps>0$, set
  \begin{equation}
    \gamma(u,\eps)
    :=\lim_{t\to\infty}\frac1t
    \log\norm{\e^{tL_{\eps,u}}}_{\Lp^2_\sigma\to\Lp^2_\sigma}.
    \label{eq:main:semigroup-growth-rate}
  \end{equation}
  The limit exists because $t\mapsto\log\norm{\e^{tL_{\eps,u}}}$ is subadditive
  and every $\C_0$-semigroup is locally bounded in time
  [Engel--Nagel, Prop.~IV.2.2].
\end{definition}

The Lean counterpart is \texttt{FastDynamo.gamma}
(\texttt{FastDynamo/Dynamics/Assembly/MainGrowthRate.lean}): the growth bound
(Definition~\ref{def:dynamics:growth-bound}) of
\texttt{literalInductionSemigroup} on the closed solenoidal subspace
\texttt{literalSolenoidalSubspace} of $\Lp^2(\T^3;\CC^3)$.  The semigroup is a
$\C_0$-semigroup, constructed and not assumed
(Theorem~\ref{thm:dynamics:solenoidal-semigroup}).  Three points of the encoding.

\begin{itemize}
  \item \emph{Which semigroup.}  \texttt{literalInductionSemigroup} is the isometric transport
    to the literal three-torus of \texttt{solenoidalInductionSemigroup}, whose generator is the
    paper's $L_{\eps,u}$ with domain $\Hs^2\cap\Lp^2_\sigma$, in the equivalent form recorded at
    Theorem~\ref{thm:dynamics:solenoidal-generator}: the values are the auxiliary action
    $\mathcal A_\eps B$ of Definition~\ref{def:main:induction-operator}, which on
    divergence-free $\Hs^2$ fields is the curl form computed in $\Hs^{-2}$.  A
    $\C_0$-semigroup is determined by its generator
    (\texttt{eq\_solenoidalInductionSemigroup\_of\_generator\_eq},
    \texttt{gamma\_eq\_growthBound\_of\_generator\_eq}), so \texttt{gamma} is the growth bound
    of the paper's $\e^{tL_{\eps,u}}$.  The route to the semigroup differs from the paper's:
    Lean generates on all of $\Lp^2(\T^3;\CC^3)$ with the auxiliary operator, restricts to
    $\Lp^2_\sigma$ and identifies the generator afterwards, where the paper generates on
    $\Lp^2_\sigma$ by perturbation of $\eps\Delta$ (Lemma~A.1 of the paper).
  \item \emph{$\limsup$ and $\lim$.}  Lean's \texttt{growthBound} is defined as a
    $\limsup$ in $[-\infty,\infty]$, so that it is defined for every semigroup with no
    preliminary work; Lemma~\ref{lem:dynamics:growth-bound-limit} proves that it is the limit
    for every $\C_0$-semigroup on a Banach space, by the paper's argument, and
    Lemma~\ref{lem:main:growth-rate-limit} records this for $\gamma(u,\eps)$, which is
    moreover a real number in $[0,M_u]$ for every $(u,\eps)$.
  \item \emph{Model.}  The torus is the product model $\T^2\times\T$, transported
    isometrically to the literal three-torus, and the velocity is axially independent;
    \texttt{literalSolenoidalSubspace} is the transported Fourier-defined solenoidal subspace
    of the product model.  The analytic and spectral properties of the semigroup are
    Lemma~\ref{lem:semigroup-spectral-growth}.
\end{itemize}

\begin{definition}[Dynamo, fast dynamo, slow dynamo]
  \label{def:main:dynamo}
  \lean{FastDynamo.gammaOf, FastDynamo.gammaOf_eq_gamma, FastDynamo.IsDynamoAt,
    FastDynamo.IsFastDynamo, FastDynamo.IsSlowDynamo}
  \leanok
  \uses{def:main:growth-rate}
  Let $u\in\W^{1,\infty}(\T^3;\R^3)$ be time-independent and divergence-free.  We call $u$ a
  \emph{dynamo at diffusivity $\eps$} when $\gamma(u,\eps)>0$.  It is a \emph{fast dynamo} if
  \[
    \liminf_{\eps\downarrow0}\gamma(u,\eps)>0.
  \]
  If $\gamma(u,\eps)>0$ for every sufficiently small $\eps$, but $\gamma(u,\eps)\to0$, we call
  it a \emph{slow dynamo}.
\end{definition}

These are the definitions of §1.1 of the paper, verbatim in Lean
(\texttt{FastDynamo/Dynamics/Assembly/FastDynamo.lean}) for a divergence-free velocity of
the axially independent class: \texttt{IsDynamoAt} is $0<\gamma(u,\eps)$, \texttt{IsFastDynamo}
is $0<\liminf$ along $\eps\to0^+$, and \texttt{IsSlowDynamo} is ``eventually positive along
$\eps\to0^+$ and tending to $0$''.  The limits in $\eps$ are taken of \texttt{gammaOf}, the
extension of $\gamma(u,\cdot)$ to all real $\eps$ (by $0$ for $\eps\leq0$), because
\texttt{gamma} takes the proof of $0<\eps$ as an argument; along $\eps\to0^+$ only positive
$\eps$ matter.  A fast dynamo is not a slow dynamo (\texttt{IsFastDynamo.not\_isSlowDynamo}),
which the paper leaves implicit; it is recorded with
Corollary~\ref{cor:main:fast-dynamo}.

\begin{lemma}[The growth rate is a limit]
  \label{lem:main:growth-rate-limit}
  \lean{FastDynamo.gamma_tendsto, FastDynamo.gamma_ne_bot_iff, FastDynamo.gamma_tendsto_real,
    FastDynamo.gamma_eq_lim, FastDynamo.zero_le_gamma,
    FastDynamo.mainGrowthFloor_le_gamma_lim}
  \leanok
  \uses{def:main:growth-rate, lem:dynamics:growth-bound-limit,
    lem:dynamics:mean-conservation, thm:main-growth-rate, def:main:witnesses}
  For every divergence-free $u$ and every $\eps>0$ the limit
  \eqref{eq:main:semigroup-growth-rate} exists and is a real number: there is $L\in\R$ with
  \[
    \frac1t\log\norm{\e^{tL_{\eps,u}}}_{\Lp^2_\sigma\to\Lp^2_\sigma}
    \xrightarrow[t\to\infty]{}L,
    \qquad
    L=\gamma(u,\eps),
    \qquad
    0\leq L\leq M_u .
  \]
  At the explicit witnesses of Definition~\ref{def:main:witnesses} and for every
  $0<\eps\leq\eps_0$, moreover $\gamma_0\leq L$.
\end{lemma}

This is the sentence after \texttt{eq:semigroup-growth-rate} in §1.1 of the paper (the limit
exists), together with the facts that it is finite and non-negative, which the paper does not
state.  In Lean: \texttt{gamma\_tendsto} (convergence in $[-\infty,\infty)$, from
Lemma~\ref{lem:dynamics:growth-bound-limit}), \texttt{gamma\_ne\_bot\_iff} and
\texttt{gamma\_tendsto\_real} (finiteness exactly under an exponential minorant, and then
convergence in $\R$), \texttt{gamma\_eq\_lim} and \texttt{zero\_le\_gamma} (the minorant
$\norm{\e^{tL_{\eps,u}}}\geq1$ of Lemma~\ref{lem:dynamics:mean-conservation}), and
\texttt{mainGrowthFloor\_le\_gamma\_lim} for the witnesses
(\texttt{FastDynamo/Dynamics/Assembly/MainGrowthRateLimit.lean} and
\texttt{FastDynamo/Dynamics/Assembly/InductionMeanConservation.lean}).  The convergence takes
place in \texttt{EReal}, with $\log0$ read as $-\infty$, which is the paper's $[-\infty,\infty)$
reading; the real-valued form is \texttt{gamma\_eq\_lim}.

\begin{proof}
  \leanok
  \uses{lem:dynamics:growth-bound-limit, lem:dynamics:mean-conservation,
    thm:main-growth-rate}
  Lemma~\ref{lem:dynamics:growth-bound-limit} applies to the semigroup on the closed subspace
  $\Lp^2_\sigma$ of a Hilbert space.  The minorant $\norm{\e^{tL_{\eps,u}}}\geq1$ makes the
  limit finite and non-negative, and the Hille--Yosida bound
  $\norm{\e^{tL_{\eps,u}}}\leq\e^{M_ut}$ bounds it by $M_u$.  At the witnesses,
  Theorem~\ref{thm:main-growth-rate} gives $\gamma_0\leq\gamma(u,\eps)$.
\end{proof}


\begin{definition}[The statement layer]
  \label{def:main:physical-clauses}
  \uses{def:main:induction-operator}
  \lean{FastDynamo.MainReducedClauses, FastDynamo.MainReducedTarget,
    FastDynamo.MainPhysicalClausesAt,
    FastDynamo.MainPhysicalClauses, FastDynamo.MainPhysicalConclusion,
    FastDynamo.mainPhysicalClauses_of_mainReducedClauses,
    FastDynamo.mainPhysicalConclusion_of_mainReducedTarget}
  \leanok
  \emph{The reduced clauses} $\mathrm{MRC}(u,\eps_0,\gamma_0)$ say: for every
  $0<\eps\leq\eps_0$ there are $K\in\Z\setminus\{0\}$, $\lambda\in\CC$ with
  $\re\lambda\geq\gamma_0$ and $b\in\Hs^2(\T^2;\CC^3)$ with
  $A_{\eps,K}b=\lambda b$ and $b\neq0$ in $\Lp^2$.  \emph{The reduced target}
  is
  \[
    \exists\,u\in\W^{1,\infty},\ \Div u=0,\ u\neq0,\ \
    \exists\,\eps_0,\gamma_0>0:\ \mathrm{MRC}(u,\eps_0,\gamma_0).
  \]
  \emph{The physical clauses} $\mathrm{MPC}(u,\eps_0,\gamma_0)$ are the
  sixteen conjuncts of the spectral half of Theorem~\ref{thm:main} at a fixed
  velocity and fixed constants, and \emph{the physical conclusion} is the same
  outer existential over $u,\eps_0,\gamma_0$ with $\mathrm{MPC}$ in place of
  $\mathrm{MRC}$.  The reduction
  \[
    \mathrm{MRC}(u,\eps_0,\gamma_0)\ \Longrightarrow\ \mathrm{MPC}(u,\eps_0,\gamma_0)
    \qquad(\Div u=0,\ \gamma_0>0)
  \]
  is the content of the finish layer.
\end{definition}

\begin{proof}
  \leanok
  \uses{def:finish:real-solution, prop:finish:real-strong-mode,
    lem:positive-eigenfunction-solenoidal, lem:real-axial-mode}
  Given a reduced eigenpair, the real strong-mode certificate of the finish
  layer supplies every physical clause at once: the lift
  $V=\e^{2\pi iKx_3}b$ is an eigenmode of $\mathcal{A}_\eps$ by the exact
  sector identity, it is divergence-free by
  Lemma~\ref{lem:positive-eigenfunction-solenoidal} (which needs
  $\re\lambda>0$ and $\Div u=0$) and hence an eigenmode of $L_{\eps,u}$, and
  Lemma~\ref{lem:real-axial-mode} turns the non-zero axial index into the
  non-vanishing of $\re V$ and the exact norm law for the real curve.
\end{proof}

The sixteen clauses of $\mathrm{MPC}(u,\eps_0,\gamma_0)$, in the order in
which \texttt{FastDynamo.MainPhysicalClauses} lists them, are the following.
Clauses (1)--(3) are the paper's
$K_\eps\in\Z\setminus\{0\}$, $\re\lambda_\eps\geq\gamma_0$ and
$0\neq b^\eps\in\Hs^2(\T^2;\CC^3)$.  Clause (4) is the reduced eigenpair
$A_{\eps,K}b=\lambda b$ of \eqref{eq:main:torus-sector}, kept as the reference
datum.  Clauses (5)--(7) carry the display
\eqref{eq:main-spectral-conclusion} of Theorem~\ref{thm:main}: the lift
$V^\eps=\e^{2\pi iK_\eps x_3}b^\eps$ lies in $\Hs^2(\T^3;\CC^3)$ by typing and
satisfies the auxiliary-operator equation $\mathcal{A}_\eps V^\eps=\lambda_\eps V^\eps$
(5), which the paper does not state, it is
divergence-free, hence lies in $\Lp^2_\sigma(\T^3)$ (6), and it satisfies the
curl-form equation $L_{\eps,u}V^\eps=\lambda_\eps V^\eps$ (7).  Clause (8) is
$B_{\rm in}^\eps=\re V^\eps\neq0$.  Clauses (9)--(11) and (15) are the
real-solution law of Theorem~\ref{thm:main}: the explicit real $\Lp^2$ curve
$B^\eps(t)=\re(\e^{t\lambda_\eps}V^\eps)$ never vanishes (9), obeys
$\norm{B^\eps(t)}_2=\e^{t\re\lambda_\eps}\norm{B_{\rm in}^\eps}_2$ (10) and the
uniform bound $\e^{t\gamma_0}\norm{B_{\rm in}^\eps}_2\leq\norm{B^\eps(t)}_2$
for $t\geq0$ (11), and is the real field represented by the $\Hs^2$ state
curve (15).  Clauses (12)--(14) and (16) say that this $\Hs^2$ state curve
starts at $\re V^\eps$, is fixed by conjugation (that is, real) at every time,
is divergence-free at every time, and is a strong solution of the induction
equation.

Two differences from the paper are worth recording.  The norm law (10) is
proved for \emph{every} real time, not only for $t\geq0$ as in
\eqref{eq:main-thm-growth}, and the uniform bound (11) with the
$\eps$-independent rate $\gamma_0$ is an extra clause, an immediate consequence of
\eqref{eq:main-thm-growth} and $\re\lambda_\eps\geq\gamma_0$ that the paper does not
state.  Clause (16) asserts that $B^\eps$ is a strong solution; that it is \emph{the}
solution, as the paper's ``its solution'' says, is
Lemma~\ref{lem:dynamics:physical-uniqueness}: forward strong solutions of the curl-form
equation in $\Hs^{-2}$ are unique.  The semigroup on $\Lp^2_\sigma(\T^3)$ carries the datum
$B_{\rm in}^\eps$ along exactly the curve $B^\eps$, and every forward strong solution with
datum $B_{\rm in}^\eps$ is that orbit; this is \texttt{FastDynamo.MainSemigroupOrbitClauses},
the field \texttt{semigroupOrbit} of Theorem~\ref{thm:main-strengthened}, which also asserts
the sixteen clauses of $\mathrm{MPC}$ at its own triple $(K,\lambda,b)$, so that the eigenmode,
the real datum and its solution belong to one triple.  That clause cannot
live in $\mathrm{MPC}$ itself: \texttt{FastDynamo/Main.lean} is in the import closure of the
semigroup construction, so naming the semigroup there would close a cycle.  The Lean modules
are \texttt{FastDynamo/Main.lean} (the four predicates and the reduction; the body of
$\mathrm{MPC}$ at one triple $(K,\lambda,b)$ is \texttt{MainPhysicalClausesAt}),
\texttt{FastDynamo/Finish/RealSolution.lean} (the certificate that discharges it) and
\texttt{FastDynamo/Dynamics/Assembly/SemigroupOrbit.lean} (the orbit clause).

\section{The main theorem}

\begin{theorem}[A fast dynamo with zero topological entropy]
  \label{thm:main}
  \lean{FastDynamo.MainTheoremComplete, FastDynamo.mainTheoremComplete_holds}
  \leanok
  \uses{def:main:induction-operator, def:main:growth-rate, thm:main-growth-rate,
    prop:particle-flow, cor:ideal-growth}
  There exist a real-valued autonomous velocity field
  \[
    u\in\W^{1,\infty}(\T^3;\R^3),\qquad\Div u=0,
  \]
  and constants $\eps_0>0$, $\gamma_0>0$, such that
  $\gamma(u,\eps)\geq\gamma_0$ for every $0<\eps\leq\eps_0$.  In particular,
  \begin{equation}
    \liminf_{\eps\downarrow0}\gamma(u,\eps)\geq\gamma_0.
    \label{eq:main-growth-conclusion}
  \end{equation}
  More precisely, for every $0<\eps\leq\eps_0$ there are
  $K_\eps\in\Z\setminus\{0\}$, $\lambda_\eps\in\CC$ and
  $0\neq b^\eps\in\Hs^2(\T^2;\CC^3)$ such that the complex magnetic eigenmode
  \begin{equation}
    V^\eps(y,x_3):=\e^{2\pi iK_\eps x_3}b^\eps(y)
    \in\Hs^2(\T^3;\CC^3)\cap\Lp^2_\sigma(\T^3),
    \quad
    L_{\eps,u}V^\eps=\lambda_\eps V^\eps,
    \quad
    \re\lambda_\eps\geq\gamma_0.
    \label{eq:main-spectral-conclusion}
  \end{equation}
  Moreover the real-valued datum
  $B_{\rm in}^\eps:=\re V^\eps\in\Hs^2(\T^3;\R^3)\cap\Lp^2_\sigma(\T^3)$ is
  non-zero, and its solution satisfies
  \begin{equation}
    \norm{B^\eps(t)}_{\Lp^2(\T^3)}
    =\e^{t\re\lambda_\eps}\norm{B_{\rm in}^\eps}_{\Lp^2(\T^3)}
    \qquad\text{for every }t\geq0.
    \label{eq:main-thm-growth}
  \end{equation}

  Moreover, $h_{\mathrm{top}}(\Phi_t)=0$ for all $t\in\R$ and
  \[
    \lim_{t\to\infty}\frac1t
    \log\norm{\mathcal T_0(t)}_{\Lp^2_\sigma\to\Lp^2_\sigma}=0,
  \]
  where $\Phi_t$ is the particle flow of $u$
  (Definition~\ref{def:dynamics:particle-flow}) and $\mathcal T_0(t)$ the solution
  group on $\Lp^2_\sigma(\T^3)$ of the ideal induction equation
  (Definition~\ref{def:dynamics:ideal-induction-group}).
\end{theorem}

\begin{proof}
  \leanok
  \uses{thm:main-spectral, thm:main-growth-rate, lem:main:growth-rate-limit,
    lem:real-axial-mode, prop:particle-flow, cor:ideal-growth, thm:main-strengthened,
    lem:dynamics:physical-uniqueness}
  Everything except the growth clause and \eqref{eq:main-growth-conclusion} is
  Theorem~\ref{thm:main-spectral}; the dynamical half is
  Proposition~\ref{prop:particle-flow} and Corollary~\ref{cor:ideal-growth}, both of which
  are formalised (Proposition~\ref{prop:main-lagrangian}).  The growth clause is
  Theorem~\ref{thm:main-growth-rate}: the eigenvalue $\lambda_\eps$ produced
  for each $0<\eps\leq\eps_0$ lies in the state space $\Lp^2_\sigma(\T^3)$ and
  the semigroup orbit of its eigenmode grows exactly like
  $\e^{t\re\lambda_\eps}$, which gives
  $\gamma(u,\eps)\geq\re\lambda_\eps\geq\gamma_0$; taking the lower limit as
  $\eps\downarrow0$ finishes.

  The paper's own route to the growth clause is the same in shape.  It reads the
  exact growth law \eqref{eq:main-thm-growth} of the \emph{real} datum as an exponential
  minorant of the operator norm,
  $\norm{\e^{tL_{\eps,u}}}\geq\norm{B^\eps(t)}_2/\norm{B_{\rm in}^\eps}_2
  =\e^{t\re\lambda_\eps}$, and concludes
  $\gamma(u,\eps)\geq\re\lambda_\eps\geq\gamma_0$; the semigroup appendix,
  Lemma~\ref{lem:semigroup-spectral-growth}, is not invoked here at all.  Lean divides by the
  norm of the \emph{complex} eigenmode $V^\eps$ rather than of the real datum
  $B_{\rm in}^\eps$ --- the two give the same inequality, $V^\eps\neq0$ being part of the
  definition of an eigenpair --- and \eqref{eq:main-thm-growth} itself, which the paper
  obtains from Lemma~\ref{lem:real-axial-mode}, is clause (10) of
  Definition~\ref{def:main:physical-clauses}.
\end{proof}

The Lean statement is \texttt{FastDynamo.MainTheoremComplete}, proved by
\texttt{mainTheoremComplete\_holds} (\texttt{FastDynamo/MainTheoremStrengthened.lean},
\texttt{FastDynamo/Dynamics/MainLagrangianClauses.lean}); it is
Theorem~\ref{thm:main-complete}, whose fields are the strengthened spectral statement
(Theorem~\ref{thm:main-strengthened}), the two dynamical clauses
(\texttt{MainTheoremDynamical}) and the larger Lagrangian bundle
(Proposition~\ref{prop:main-lagrangian}).  Clause by clause: the growth clause and
\eqref{eq:main-growth-conclusion} are Theorem~\ref{thm:main-growth-rate}
(\texttt{gamma\_ge\_gamma0}, \texttt{mainGrowthFloor\_le\_liminf\_gammaOf}), with
$\gamma(u,\eps)$ the limit of Lemma~\ref{lem:main:growth-rate-limit};
\eqref{eq:main-spectral-conclusion}, the non-vanishing of $B_{\rm in}^\eps$ and the exact norm
law \eqref{eq:main-thm-growth} are Theorem~\ref{thm:main-spectral}; ``its solution'' is the
semigroup-orbit clause of Theorem~\ref{thm:main-strengthened}, which says that $B_{\rm in}^\eps$
lies in $\Lp^2_\sigma$, that its orbit under $\e^{tL_{\eps,u}}$ is the explicit real curve
$B^\eps(t)=\re(\e^{t\lambda_\eps}V^\eps)$ for $t\geq0$, and that \emph{every} forward strong
solution of the curl-form equation with datum $B_{\rm in}^\eps$ is that orbit and obeys
\eqref{eq:main-thm-growth} (Lemma~\ref{lem:dynamics:physical-uniqueness}); and the two
dynamical clauses are \texttt{FastDynamo.mainTheoremDynamical\_holds}, with $\mathcal T_0$ the
unique solution group of the ideal equation (Corollary~\ref{cor:ideal-growth}).

The Lean statement differs from the paper's in these respects, each recorded at the node
named.
\begin{itemize}
  \item \emph{Shape of the statement.}  The clauses are fields and theorems at the explicit
    witnesses of Definition~\ref{def:main:witnesses}, not one proposition of the form
    $\exists u,\eps_0,\gamma_0$, and the per-diffusivity predicates
    \texttt{MainPhysicalClauses}, \texttt{MainSemigroupOrbitClauses} and
    \texttt{MainLiteralTorusClauses} each carry their own existential over $(K,\lambda,b)$.
    \texttt{MainSemigroupOrbitClauses} contains the sixteen clauses of
    Definition~\ref{def:main:physical-clauses} (\texttt{MainPhysicalClausesAt}) at its own
    triple, so the eigenmode, the real datum, the strong solution, the semigroup orbit and the
    uniqueness clause hold at one $(K,\lambda,b)$.  The triple of
    \texttt{MainLiteralTorusClauses} is the one of \texttt{mainSemigroupOrbitClauses\_holds}
    by its proof, not by its statement, and the field \texttt{physicalClauses} of
    Theorem~\ref{thm:main-strengthened} repeats the spectral clauses with a separately chosen
    triple (the paragraph after Theorem~\ref{thm:main-strengthened}).
  \item \emph{Velocity class and torus.}  The velocity ranges over the axially independent
    class \texttt{TransverseRealW1Inf}, so the existential is over a smaller set than the
    paper's and is stronger; the magnetic clauses live on the product model $\T^2\times\T$,
    transported isometrically to the literal three-torus (reading guide,
    Lemma~\ref{lem:main:literal-torus}).
  \item \emph{The operator} $L_{\eps,u}$ is realised as the curl form with values in
    $\Hs^{-2}$ in the statement layer and, as generator, in the equivalent auxiliary form on
    $\Hs^2\cap\Lp^2_\sigma$ (Definition~\ref{def:main:induction-operator},
    Theorem~\ref{thm:dynamics:solenoidal-generator}).
  \item \emph{The solution class} of ``its solution'' is not named by the paper; Lean takes the
    forward strong solutions of the curl-form equation in $\Hs^{-2}$
    (Lemma~\ref{lem:dynamics:physical-uniqueness}), and the clause adds the lower bound
    $\e^{t\gamma_0}\norm{B_{\rm in}^\eps}_2\leq\norm{B^\eps(t)}_2$, which the theorem does not
    state.
  \item \emph{The ideal solution group} is unique among the solution groups of the weak form
    of Definition~\ref{def:dynamics:ideal-weak-form}, whose test fields and time integration
    differ from the paper's distributional notion as recorded there.
\end{itemize}

\textbf{The ideal--resistive growth gap.}  The paper folds this contrast into
Theorem~\ref{thm:main}, whose closing line is its left member and
whose growth clause is its right member: for the velocity constructed here the ideal
induction group has exponential growth rate zero, while every sufficiently small positive
diffusivity grows at least at the rate $\gamma_0>0$,
\[
  \lim_{t\to\infty}\frac1t
  \log\norm{\mathcal T_0(t)}_{\Lp^2_\sigma\to\Lp^2_\sigma}
  =0<\gamma_0\leq\liminf_{\eps\downarrow0}\gamma(u,\eps).
\]
The chain is not a statement of the paper, and therefore not a node of this
blueprint; the paper's own qualitative reading of it is the second paragraph of its
Remark~1.3 (\texttt{rem:subsequent-resolution}).  In Lean,
\texttt{FastDynamo.ideal\_resistive\_gap\_holds} is the display itself, and
\texttt{FastDynamo.ideal\_resistive\_gap\_at\_diffusivity\_holds} is the sharper fixed-$\eps$
packaging, with $\gamma_0\leq\gamma(u,\eps)$ at every $0<\eps\leq\eps_0$ in place of the lower
limit.  Its left end is \texttt{HasVanishingIdealOperatorGrowthRate}
(Corollary~\ref{cor:ideal-growth}), its middle is $\gamma_0>0$, and its right end is
\texttt{mainGrowthFloor\_le\_liminf\_gammaOf}, the lower limit being taken along the filter
$\mathcal N_{>0}$ of the total extension \texttt{gammaOf} of $\gamma(u,\cdot)$.

\begin{theorem}[The spectral half of Theorem~\ref{thm:main}]
  \label{thm:main-spectral}
  \uses{def:main:physical-clauses}
  \lean{FastDynamo.MainPhysicalConclusion,
    FastDynamo.mainPhysicalConclusion_holds,
    FastDynamo.main_theorem_spectral, FastDynamo.mainReducedTarget_holds}
  \leanok
  There exist a real-valued autonomous $u\in\W^{1,\infty}(\T^3;\R^3)$ with
  $\Div u=0$ and $u\neq0$, and constants $\eps_0,\gamma_0>0$, such that for
  every $0<\eps\leq\eps_0$ the sixteen clauses
  $\mathrm{MPC}(u,\eps_0,\gamma_0)$ of
  Definition~\ref{def:main:physical-clauses} hold: an axial index
  $K_\eps\neq0$, an eigenvalue with $\re\lambda_\eps\geq\gamma_0$ and a
  non-zero $b^\eps\in\Hs^2(\T^2;\CC^3)$ producing the solenoidal $\Hs^2$
  eigenmode $V^\eps=\e^{2\pi iK_\eps x_3}b^\eps$ of
  \eqref{eq:main-spectral-conclusion}, the non-zero real datum
  $B_{\rm in}^\eps=\re V^\eps$, and the explicit real, divergence-free strong
  solution $B^\eps$ with
  $\norm{B^\eps(t)}_2=\e^{t\re\lambda_\eps}\norm{B_{\rm in}^\eps}_2$ for every
  $t\in\R$ and $\norm{B^\eps(t)}_2\geq\e^{t\gamma_0}\norm{B_{\rm in}^\eps}_2$
  for $t\geq0$.

  Equivalently: Theorem~\ref{thm:main} holds with its growth clause and its two dynamical
  clauses deleted, and ``its solution'' read as ``an explicit real, divergence-free strong
  solution with that datum''.
\end{theorem}

\begin{proof}
  \leanok
  \uses{def:transfer:parameter-choice, prop:global-spectral-transfer,
    lem:global-riesz-nonzero, lem:axial-recovery,
    lem:real-axial-mode, def:finish:real-solution, prop:finish:real-strong-mode}
  Fix the global parameter chain of
  Definition~\ref{def:transfer:parameter-choice}; it pins the velocity $u$,
  the threshold $\eps_0$ and, for each $0<\eps\leq\eps_0$, a matching cell and
  its axial index $K_\eps\neq0$.  Proposition~\ref{prop:global-spectral-transfer}
  gives, at that chain and for every such $\eps$, a resolvent of
  $A_{\eps,K_\eps}$ on the selected contour whose Riesz projection is non-zero;
  this is the step whose non-vanishing is
  Lemma~\ref{lem:global-riesz-nonzero}.  The projection has finite rank, so a
  non-zero Riesz projection yields an eigenpair of $A_{\eps,K_\eps}$
  inside the contour, whose real part is at least the contour's growth floor
  $\gamma_0=g_*$ -- a closed term of the fixed local seed, independent of
  $\eps$, of the velocity, of the sector and of the placement.  The eigenpair is
  \emph{horizontal}; Lemma~\ref{lem:axial-recovery} recovers the third component and
  delivers the reduced target.  The reduction of
  Definition~\ref{def:main:physical-clauses}, which uses
  Lemma~\ref{lem:real-axial-mode} through the certificate of
  Definition~\ref{def:finish:real-solution}, then produces all sixteen physical
  clauses.
\end{proof}

The Lean proof follows this route.  \texttt{mainReducedTarget\_holds} and
\texttt{mainPhysicalConclusion\_holds} in \texttt{FastDynamo/MainTheorem.lean} are obtained
from \texttt{mainReducedClauses\_holds}
(\texttt{FastDynamo/Global/MainTheoremWitnesses.lean}), which is three lines through the
modules
\texttt{Global/GlobalParameterChoiceConcrete.lean} $\to$
\texttt{Global/GlobalSpectralTransferConcrete.lean} $\to$
\texttt{Global/TorusGlobalSpectralTransfer.lean}; the finish layer is entered exactly once,
through \texttt{realStrongModeCertificate\_of\_reducedEigenpair}.
Corollary~\ref{cor:magnetic-scale} (the two-sided bounds
$\abs{K_\eps}\asymp\eps^{-1/2}$ and $L_B^\eps\asymp\sqrt\eps$) is \emph{not} used anywhere
in this chain: it is a separate quantitative statement about the modes once they have been
produced, and no step of the Lean proof invokes it.

\begin{theorem}[The growth-rate clause of Theorem~\ref{thm:main}]
  \label{thm:main-growth-rate}
  \uses{def:main:growth-rate, thm:main-spectral, thm:dynamics:solenoidal-semigroup,
    lem:dynamics:solenoidal-subspace, def:main:witnesses}
  \lean{FastDynamo.gamma_ge_gamma0, FastDynamo.mainGrowthFloor_le_gamma_le_stretching,
    FastDynamo.gamma_le_stretching,
    FastDynamo.gamma_ne_bot, FastDynamo.gamma_ne_top,
    FastDynamo.gammaOf, FastDynamo.gammaOf_eq_gamma,
    FastDynamo.mainGrowthFloor_le_liminf_gammaOf}
  \leanok
  At the explicit witnesses $u$, $\eps_0$, $\gamma_0$ of
  Definition~\ref{def:main:witnesses}, for every $0<\eps\leq\eps_0$,
  \[
    \gamma_0\ \leq\ \gamma(u,\eps)\ \leq\ M_u,
    \qquad
    M_u=\norm{\nabla_\perp u}_\infty,
  \]
  $\gamma(u,\eps)$ being the growth bound (Definition~\ref{def:main:growth-rate})
  of the semigroup on $\Lp^2_\sigma(\T^3;\CC^3)$ of
  Theorem~\ref{thm:dynamics:solenoidal-semigroup}.  In particular
  $\gamma(u,\eps)$ is neither $-\infty$ nor $+\infty$.  Since $\gamma_0$ does not
  depend on $\eps$, this is the growth clause of Theorem~\ref{thm:main}, and with it
  \eqref{eq:main-growth-conclusion}.
\end{theorem}

The lower bound is \texttt{FastDynamo.gamma\_ge\_gamma0}, the upper bound
\texttt{gamma\_le\_stretching}, and the pair
\texttt{mainGrowthFloor\_le\_gamma\_le\_stretching}
(\texttt{FastDynamo/Dynamics/Assembly/MainGrowthRate.lean}); both are fields of
the strengthened theorem (Theorem~\ref{thm:main-strengthened}).  With
Lemma~\ref{lem:main:growth-rate-limit}, which shows that $\gamma(u,\eps)$ is the paper's limit,
this is \eqref{eq:main-growth-conclusion} in the form displayed there; the packaged
statement is \texttt{mainGrowthFloor\_le\_gamma\_lim}.  Two smaller points.  The
paper's trailing restatement
$\liminf_{\eps\downarrow0}\gamma(u,\eps)\geq\gamma_0$ is
\texttt{mainGrowthFloor\_le\_liminf\_gammaOf}, the $\liminf$ along
$\mathcal N_{>0}$ of the total extension \texttt{gammaOf} of $\gamma(u,\cdot)$ ---
\texttt{gamma} carries the proof $0<\eps$ as an argument, so a filter limit in $\eps$ needs
that extension, which agrees with $\gamma$ on $(0,\infty)$ by
\texttt{gammaOf\_eq\_gamma}.  It is weaker than the per-diffusivity bound and is recorded
because it is the form of \eqref{eq:main-growth-conclusion} and of the paper's definition of
a fast dynamo.  The interval is the paper's: \texttt{gamma\_ge\_gamma0} and
\texttt{MainTheoremStrengthened.growthRate} carry the hypotheses $0<\eps$ and
$\eps\leq\eps_0$, the closed interval $0<\eps\leq\eps_0$ on which the paper states the growth
clause.  The semigroup is the paper's $\e^{tL_{\eps,u}}$, obtained by restriction and
identified by its generator (Theorem~\ref{thm:dynamics:solenoidal-generator}).  The upper bound is not in the
paper; it is
recorded because it is what makes the sandwich informative --- without it the
lower bound would be compatible with $\gamma=+\infty$ --- and because it is
uniform in the diffusivity, which is the paper's energy competition
(\texttt{eq:energy-competition}) read at the level of the semigroup.  The floor $\gamma_0$ is
a closed term of the fixed local seed with $\gamma_0\geq\tfrac34(448-4\pi^2)\geq288$; these
numerals are choices of the formalisation, not of the paper
(Definition~\ref{def:main:witnesses}).

\begin{proof}
  \leanok
  \uses{thm:main-spectral, thm:dynamics:solenoidal-semigroup,
    lem:dynamics:solenoidal-subspace, lem:positive-eigenfunction-solenoidal}
  This is the paper's route, with one extra Lean-side obligation.  Fix
  $0<\eps\leq\eps_0$.  The reduced clauses of Theorem~\ref{thm:main-spectral}
  supply an axial index $K_\eps\neq0$, an eigenvalue $\lambda_\eps$ with
  $\re\lambda_\eps\geq\gamma_0$ and a non-zero $b^\eps\in\Hs^2(\T^2;\CC^3)$ with
  $A_{\eps,K_\eps}b^\eps=\lambda_\eps b^\eps$; the exact sector identity promotes
  this to an eigenpair of $\mathcal{A}_\eps$ at the lift
  $V^\eps=\e^{2\pi iK_\eps x_3}b^\eps$.  The extra obligation is that $V^\eps$
  lie in the state space: since $\gamma_0>0$ forces $\re\lambda_\eps>0$, the
  Lemma~\ref{lem:positive-eigenfunction-solenoidal}, in the membership
  form of Lemma~\ref{lem:dynamics:solenoidal-subspace}, puts $V^\eps$ in
  $\Lp^2_\sigma$.  Theorem~\ref{thm:dynamics:solenoidal-semigroup} then converts
  the eigenvalue into $\re\lambda_\eps\leq\gamma(u,\eps)$ and supplies the
  ceiling $\gamma(u,\eps)\leq M_u$.  Since the lower bound holds for every
  $0<\eps\leq\eps_0$ with an $\eps$-independent $\gamma_0$, the lower limit as
  $\eps\downarrow0$ is at least $\gamma_0$.
\end{proof}

\begin{corollary}[The constructed velocity is a fast dynamo]
  \label{cor:main:fast-dynamo}
  \lean{FastDynamo.isDynamoAt_mainVelocity, FastDynamo.isFastDynamo_mainVelocity,
    FastDynamo.IsFastDynamo.not_isSlowDynamo}
  \leanok
  \uses{def:main:dynamo, thm:main-growth-rate, def:main:witnesses}
  The velocity $u$ of Definition~\ref{def:main:witnesses} is a dynamo at every diffusivity
  $0<\eps\leq\eps_0$ and a fast dynamo in the sense of Definition~\ref{def:main:dynamo}; in
  particular it is not a slow dynamo.
\end{corollary}

The paper draws this conclusion in its title and abstract and through
\eqref{eq:main-growth-conclusion}, which is its definition of a fast dynamo; it does not state
it as a separate result.  Module \texttt{FastDynamo/Dynamics/Assembly/FastDynamo.lean}.

\begin{proof}
  \leanok
  \uses{thm:main-growth-rate}
  $\gamma(u,\eps)\geq\gamma_0>0$ for $0<\eps\leq\eps_0$ by Theorem~\ref{thm:main-growth-rate},
  and $\liminf_{\eps\downarrow0}\gamma(u,\eps)\geq\gamma_0>0$.  If $\gamma(u,\eps)\to0$ the
  lower limit would be $0$.
\end{proof}

\section{The strengthened statement}

The proof produces closed terms for the velocity and both constants, and it
proves several statements that are strictly stronger than
Theorem~\ref{thm:main-spectral}.  They are collected in one structure.

\begin{definition}[The explicit witnesses]
  \label{def:main:witnesses}
  \uses{def:transfer:parameter-choice, def:velocity:global-velocity,
    def:local:seed}
  \lean{FastDynamo.mainVelocity,
    FastDynamo.mainEps0,
    FastDynamo.mainGrowthFloor,
    FastDynamo.mainEps0_eq,
    FastDynamo.mainEps0_pos,
    FastDynamo.mainGrowthFloor_eq,
    FastDynamo.mainGrowthFloor_pos,
    FastDynamo.mainGrowthFloor_ge_three_quarters,
    FastDynamo.mainGrowthFloor_ge_numeric,
    FastDynamo.mainVelocity_isTransverseDivergenceFree,
    FastDynamo.mainVelocity_ne_zero,
    FastDynamo.seed_kStar, FastDynamo.seed_wavenumber}
  \leanok
  Let $\Theta$ be the global parameter chain of
  Definition~\ref{def:transfer:parameter-choice} produced by the closure, let
  $u:=u_\Theta$ be its global velocity
  (Definition~\ref{def:velocity:global-velocity}), and put
  \[
    \eps_0:=\delta_+\frac{k_*^2}{K_0^2},
    \qquad
    \gamma_0:=g_*=\re(\text{centre }\Gamma)-\text{radius }\Gamma .
  \]
  Here $\delta_+$ is the seed's upper rescaled-diffusivity endpoint, $k_*=1$ the seed's
  integer wavenumber (so $N=2\pi$), $K_0$ the base axial index of the chain and $\Gamma$ the
  seed's isolating circle, so that $g_*=\min_{\zeta\in\cG}\re\zeta$.
\end{definition}

The threshold $\eps_0=\delta_+/K_0^2$ (at $k_*=1$) is the paper's value, the display
after \texttt{eq:diffusivity-intervals} in §3.2 of the paper, with no further cap;
\texttt{mainEps0\_eq} is that formula, and $\gamma_0=g_*$ is the choice made in the paper's
proof of Theorem~1.2.  The growth floor is a closed term of the fixed local seed alone: it
depends on neither $\eps$, nor the velocity, nor the sector, nor the axial shift, nor the
placement.

The numerals that follow are choices of the formalisation, not of the paper.  The paper
takes the radial scale $\varrho$ of the segment $\kappa=\varrho(1+it)$ ``sufficiently
large'', obtains $\re\lambda_*\geq\tfrac7{16}\varrho^2-4\pi^2$
(\texttt{eq:local-eigenvalue}), and fixes the isolating circle $\Gamma$ only qualitatively,
so it gives no numerical value of $\gamma_0=g_*$.  Lean takes $\varrho=32$ and $N=2\pi$,
which gives $\re\lambda_*\geq448-4\pi^2$, and a circle centred at $\lambda_*$ of radius at
most $\re\lambda_*/4$, which gives $g_*\geq\tfrac34\re\lambda_*$
(Lemmas~\ref{lem:local:segment-root} and~\ref{lem:local:isolation}); together they give
\[
  \gamma_0\ \geq\ \tfrac34\bigl(448-4\pi^2\bigr)\ \approx\ 306.4,
  \qquad\text{and in particular}\qquad
  \gamma_0\ \geq\ 288 .
\]
In Lean these are \texttt{FastDynamo.mainEps0\_eq},
\texttt{FastDynamo.mainGrowthFloor\_eq},
\texttt{FastDynamo.mainGrowthFloor\_ge\_three\_quarters} and
\texttt{FastDynamo.mainGrowthFloor\_ge\_numeric}
(\texttt{FastDynamo/Global/MainTheoremWitnesses.lean}); the numeral $288$ comes from
$\pi\leq4$, and the sharp value is $\approx306.4$.  Note that $\eps_0$ is explicit only as a
formula: $K_0$ is obtained from an existence statement about the closure thresholds and is
not computed, so no numeral for $\eps_0$ is claimed.

\begin{lemma}[Normalisation]
  \label{lem:main:normalisation}
  \uses{def:main:induction-operator, def:main:physical-clauses,
    def:main:witnesses, def:main:growth-rate, def:dynamics:particle-flow,
    def:dynamics:ideal-induction-group}
  \lean{FastDynamo.reducedInductionAction_smul,
    FastDynamo.reducedInductionOperator_isEigenpair_smul_iff,
    FastDynamo.mainReducedClauses_smul,
    FastDynamo.MainPhysicalConclusionNormalized,
    FastDynamo.mainPhysicalConclusionNormalized_of_mainReducedTarget,
    FastDynamo.mainPhysicalClauses_normalizedVelocity,
    FastDynamo.productInductionAction_smul,
    FastDynamo.productInductionAction_normalizedVelocity,
    FastDynamo.StronglyContinuousSemigroup.timeRescale,
    FastDynamo.StronglyContinuousSemigroup.growthBound_timeRescale,
    FastDynamo.literalInductionSemigroup_normalizedVelocity,
    FastDynamo.gamma_normalizedVelocity,
    FastDynamo.normalizedGrowthFloor_le_gamma,
    FastDynamo.normalizedGrowthFloor_le_liminf_gammaOf,
    FastDynamo.mainSemigroupOrbitClauses_normalizedVelocity,
    FastDynamo.IsParticleFlow.timeRescale, FastDynamo.IsIdealInductionGroup.timeRescale,
    FastDynamo.normalizedLiteralVelocity,
    FastDynamo.isIdealInductionGroup_normalizedLiteralVelocity,
    FastDynamo.normalizedLiteralFlow_unique,
    FastDynamo.MainTheoremNormalized, FastDynamo.mainTheoremNormalized_holds}
  \leanok
  Put $M=\max\{1,\norm{u}_{\W^{1,\infty}}\}$ and $\widetilde u=u/M$.  Then
  \[
    L_{\eps,\widetilde u}=M^{-1}L_{M\eps,u},
    \qquad
    \gamma(\widetilde u,\eps)=M^{-1}\gamma(u,M\eps).
  \]
  Hence Theorem~\ref{thm:main} remains valid for $\widetilde u$, with threshold
  $\widetilde\eps=\eps_0/M$ and lower bound $\widetilde\gamma=\gamma_0/M$ on the growth rate.
  Moreover, the particle flow and ideal induction group for $\widetilde u$ are the
  corresponding objects for $u$ evaluated at time $t/M$.  This preserves the ideal growth rate
  and the vanishing of topological entropy.  In particular, the velocity may be required to
  satisfy $\norm{\widetilde u}_{\W^{1,\infty}}\leq1$ without changing any of the conclusions of
  Theorem~\ref{thm:main}.
\end{lemma}

This is Remark~1.6 of the paper (unlabelled), whose wording the statement follows.  In Lean it
is the structure \texttt{MainTheoremNormalized} at $\widetilde u=$
\texttt{normalizedVelocity mainVelocity}, proved by \texttt{mainTheoremNormalized\_holds}
(\texttt{FastDynamo/MainTheoremNormalized.lean}): $\norm{\widetilde u}\leq1$; the two
identities (\texttt{action\_eq}, and \texttt{gamma\_eq} from the identity of semigroups
$\e^{tL_{\eps,\widetilde u}}=\e^{(t/M)L_{M\eps,u}}$, \texttt{semigroup\_eq}); every clause of
Theorem~\ref{thm:main} for $\widetilde u$ with $M^{-1}\eps_0$ and $M^{-1}\gamma_0$ (the spectral
clauses, the growth-rate clause and its $\liminf$ form, the semigroup-orbit and solution
clause, and the two dynamical clauses); the flow $t\mapsto\Phi_{t/M}$, which is the only
particle flow of $\widetilde u$, and the ideal solution group $t\mapsto\mathcal T_0(t/M)$; and
Proposition~\ref{prop:main-lagrangian} for $\widetilde u$.  The deviations:
\begin{itemize}
  \item The first identity is stated for the auxiliary all-frequency action $\mathcal A_\eps$
    on all of $\Hs^2$ (\texttt{productInductionAction\_normalizedVelocity}), the action of the
    generator, which on divergence-free fields is the curl form
    (Theorem~\ref{thm:dynamics:solenoidal-generator}); for the reduced operator the same
    homogeneity $A_{c\eps,cu,K}=c\,A_{\eps,u,K}$ is \texttt{reducedInductionAction\_smul}.
  \item The norm in $M$ is the norm of \texttt{TransverseRealW1Inf},
    $\norm u=\max(\norm u_\infty,\norm{\partial_1u}_\infty,\norm{\partial_2u}_\infty)$ with
    $\norm\cdot_\infty$ the essential supremum of the Euclidean norm; the paper does not
    specify the norm $\norm\cdot_{\W^{1,\infty}}$, and every identity holds for any $M>0$.
  \item Proposition~\ref{prop:main-lagrangian} for $\widetilde u$ is proved with the constant
    $(1+M^{-1})C_{\rm L}$; the paper does not restate the constants.
\end{itemize}
The normalised witness is \emph{named}, not merely some velocity in the unit ball.  Modules
\texttt{FastDynamo/Dynamics/SemigroupTimeRescaling.lean},
\texttt{FastDynamo/Dynamics/Assembly/NormalizedGrowthRate.lean},
\texttt{FastDynamo/Dynamics/NormalizedLagrangianClauses.lean} and
\texttt{FastDynamo/Operator/ReducedInductionScaling.lean}.

\begin{proof}
  \leanok
  \uses{thm:main-spectral, thm:main-growth-rate, prop:particle-flow, cor:ideal-growth,
    lem:dynamics:semigroup-unique, lem:dynamics:physical-uniqueness}
  The three velocity-dependent terms of the induction operator are $\R$-linear in $u$ and the
  diffusion term is linear in $\eps$, which gives $L_{c\eps,cu}=c\,L_{\eps,u}$.  The semigroup
  of $cu$ at $c\eps$ and the semigroup of $u$ at $\eps$ run at time $ct$ are realised by
  operators that differ by the factor $c$, so they coincide by uniqueness of the realising
  semigroup; the growth bound of $t\mapsto S(ct)$ is $c\,\omega_0(S)$.  For the clauses, given
  $0<\eps\leq M^{-1}\eps_0$ apply the unnormalised clauses at $M\eps\leq\eps_0$ and rescale
  the eigenvalue.  If $\Phi$ is the flow of $u$, then $t\mapsto\Phi_{t/M}$ solves the flow
  equation of $u/M$, and the Cauchy formula, the weak form and uniqueness rescale with it; the
  entropy of $\Phi_{t/M}$ is zero and the logarithmic rates are unchanged by a linear change of
  time.
\end{proof}

\begin{lemma}[Qualitative regularity of the velocity]
  \label{lem:main:velocity-regularity}
  \uses{def:main:witnesses, def:velocity:global-velocity,
    def:velocity:scale-family}
  \lean{FastDynamo.ScaleFamily.TorusCellPlacement.FullVelocityConcrete.GlobalVelocityQualitative,
    FastDynamo.mainVelocityChartRepresentative_qualitative,
    FastDynamo.differentiable_mainVelocityPeriodizedRepresentative,
    FastDynamo.contDiffAt_mainVelocityPeriodizedRepresentative,
    FastDynamo.fderiv_mainVelocityPeriodizedRepresentative_accumulationPoint,
    FastDynamo.not_continuousAt_fderiv_mainVelocityPeriodizedRepresentative,
    FastDynamo.lipschitzWith_mainVelocityPeriodizedRepresentative,
    FastDynamo.lipschitzWith_mainVelocityTorusRepresentative,
    FastDynamo.mainVelocityTorusRepresentative_transverseTorusPoint_eq_periodized}
  \leanok
  Let $u$ be the velocity of Definition~\ref{def:main:witnesses} and let
  $U\colon\R^2\to\R^3$ be its $\Z^2$-periodic lift, so that $u$ is $U$ read on
  $\T^2$ and, being independent of $x_3$, on $\T^3$.  Write
  $p_\infty=(0,\tfrac12)$ for the accumulation point.  Then:
  \begin{itemize}
    \item $U$ is Lipschitz, and so are $u$ on $\T^2$ for the quotient metric
      and its axially independent extension to $\T^3$;
    \item $U$ is differentiable at every point of $\R^2$;
    \item $U$ is smooth of every order at every point outside
      $p_\infty+\Z^2$;
    \item $DU(p_\infty)=0$, yet $DU$ is discontinuous at $p_\infty$.
  \end{itemize}
  Thus $u$ is differentiable everywhere and smooth away from the accumulation
  circle $\{p_\infty\}\times\T$, its derivative vanishes on that circle, and
  $u\notin\C^1$.
\end{lemma}

\begin{proof}
  \leanok
  The local profile $v$ is smooth, and regularity is lost only because the
  cell sum has non-vanishing gradients at infinitely many shrinking horizontal
  scales.  The rescaled cells have amplitude $O(\ell_n)$ and gradients $O(1)$,
  so the disjoint sum is Lipschitz and its derivative tends to $0$ at
  $p_\infty$, which gives differentiability with $Du(p_\infty,\cdot)=0$.
  Both cut-offs are constant near each cell centre, so there the amplitude
  $\ell_n$ of the local copy cancels its dilation exactly and
  \[
    Du(p_n,x_3)=
    \begin{pmatrix}
      0      & -\Omega & 0 \\
      \Omega & 0       & 0 \\
      0      & 0       & 0
    \end{pmatrix}
    \neq0
    \qquad(x_3\in\T),
  \]
  one and the same non-zero matrix at every $n$.  Since $p_n\to p_\infty$ while
  $Du(p_\infty,\cdot)=0$, the derivative $Du$ is discontinuous at every point of the
  accumulation circle; a continuous representative of the weak gradient would have to agree
  with $Du$ everywhere, so none exists and $u\notin\C^1$.
\end{proof}

The lift $U$ is the Lean declaration
\texttt{mainVelocityPeriodizedRepresentative}, the $\Z^2$-periodisation of the
chart-side representative; the theorem
\texttt{mainVelocityTorusRepresentative\_transverseTorusPoint\_eq\_periodized}
identifies it with the pullback of the torus-side representative along the
quotient map \emph{at every point of the plane}, including the accumulation
point, which the placement chart itself misses.  The four bullets are Lean
theorems verbatim, the accumulation circle appearing as the plane point
\texttt{accumulationPoint} $=(0,\tfrac12)$ together with its integer
translates.  The proof is the paper's (§3.4 of the paper), along the cell centres:
\texttt{fderiv\_\allowbreak parameterizedFullGlobalEuclideanRealSource\_\allowbreak cellCenter\_\allowbreak eq} gives one and the
same derivative at every centre,
\texttt{lineDeriv\_parameterizedSmoothedVelocity\_zero\_apply\_one} is the entry
$(\partial_{e_1}v)(0)(e_2)=\Omega$ that makes it non-zero, and
\texttt{not\_\allowbreak continuousAt\_\allowbreak fderiv\_\allowbreak parameterizedFullGlobalEuclideanRealSource} concludes along
\texttt{tendsto\_cellCenter\_accumulationPoint}; for the lift $U$ this is
\texttt{not\_\allowbreak continuousAt\_\allowbreak fderiv\_\allowbreak mainVelocityPeriodizedRepresentative}.  The matrix itself is not assembled in
Lean; only the one entry is, which is all the non-vanishing needs.  Two caveats.  First, the final sentence of the lemma -- the
passage from the lift to a statement about a function on $\T^3$ -- is
\emph{not} itself formalised for differentiability, smoothness or the failure
of $\C^1$: the Lean development contains no \texttt{ContDiff} statement about a
map defined on $\T^3$, only the standard chart form above.  Lipschitz
continuity is the exception: it is proved torus-side as well, by
\texttt{lipschitzWith\_mainVelocityTorusRepresentative} for the quotient metric
of $\T^2$ and by \texttt{lipschitzWith\_mainLiteralVelocity} on the literal
three-torus, and the torus-side representative is almost everywhere the
$\Lp^\infty$ value of $u$.  Second, the Lipschitz constant is explicit but not
claimed optimal: it is twice the chart-side bound
\texttt{mainVelocityLipschitzBound}.  The modules are
\texttt{FastDynamo/Global/TorusGlobalVelocityRegularity.lean} and
\texttt{FastDynamo/Global/TorusGlobalVelocityLipschitz.lean}.

\begin{lemma}[The literal three-torus]
  \label{lem:main:literal-torus}
  \uses{def:main:witnesses, def:main:physical-clauses, def:finish:real-solution}
  \lean{FastDynamo.toLiteralR3L2, FastDynamo.literalRealTrajectory,
    FastDynamo.literalVelocity, FastDynamo.mainLiteralVelocity,
    FastDynamo.lipschitzWith_mainLiteralVelocity}
  \leanok
  The magnetic clauses of Theorem~\ref{thm:main-spectral} and the Lipschitz
  bound on the velocity hold on the literal three-torus
  $\T^3=(\R/\Z)^3$: for every $0<\eps\leq\eps_0$ the transported real
  trajectory $B^\eps$ lies in $\Lp^2(\T^3;\R^3)$, never vanishes, obeys
  $\norm{B^\eps(t)}_2=\e^{t\re\lambda_\eps}\norm{B^\eps(0)}_2$ for every
  $t\in\R$ and $\norm{B^\eps(t)}_2\geq\e^{t\gamma_0}\norm{B^\eps(0)}_2$ for
  $t\geq0$, and $u$ is an axially independent Lipschitz function on $\T^3$.
\end{lemma}

\begin{proof}
  \leanok
  The development works in the product model $\T^2\times\T$ with the product
  of the transverse volume and the Haar measure of the circle.  That model is
  identified with the literal $\T^3=(\R/\Z)^3$, realised in Lean as
  \texttt{Torus 3}, carrying its own volume, by a
  measure-preserving equivalence, and the induced map on $\Lp^2$ is a linear
  isometry.  Transporting the real trajectory along it preserves every norm,
  so the exact law, the non-vanishing and the growth bound transfer verbatim;
  the axially independent extension of a Lipschitz transverse field is
  Lipschitz for the sup metric of $\T^3$.
\end{proof}

\begin{lemma}[The real solution on the literal three-torus]
  \label{lem:main:literal-torus-solution}
  \uses{lem:main:literal-torus, def:main:growth-rate, lem:dynamics:physical-uniqueness}
  \lean{FastDynamo.complexifyLiteralR3L2, FastDynamo.toLiteralC3L2_complexifyProductR3L2,
    FastDynamo.MainLiteralTorusClauses,
    FastDynamo.mainLiteralTorusClauses_holds}
  \leanok
  For every $0<\eps\leq\eps_0$ there are $K\neq0$, $\lambda$ with $\re\lambda\geq\gamma_0$ and
  $b\neq0$ such that the transported real trajectory $B^\eps$ of
  Lemma~\ref{lem:main:literal-torus} obeys the norm law, the non-vanishing and the growth
  bound stated there, and there is a datum $B_{\rm in}^\eps\in\Lp^2_\sigma(\T^3;\CC^3)$, the
  complexification of $B^\eps(0)$, whose orbit under the induction semigroup
  $\e^{tL_{\eps,u}}$ on $\Lp^2_\sigma(\T^3;\CC^3)$ (the semigroup of
  Definition~\ref{def:main:growth-rate}) is the complexification of $B^\eps(t)$ for every
  $t\geq0$.
\end{lemma}

\begin{proof}
  \leanok
  The triple and the product-model datum are those of the orbit clause
  \texttt{MainSemigroupOrbitClauses} at the witnesses
  (\texttt{mainSemigroupOrbitClauses\_holds}, from
  \texttt{FastDynamo/Dynamics/Assembly/SemigroupOrbit.lean}), where the orbit of $B_{\rm in}^\eps$ under the semigroup on $\Lp^2_\sigma(\T^2\times\T)$ is the
  complexified real curve.  The semigroup on $\Lp^2_\sigma(\T^3;\CC^3)$ is that semigroup
  conjugated by the isometric transport, and the transport commutes with complexification
  (\texttt{toLiteralC3L2\_complexifyProductR3L2}: both are composition with a pointwise map,
  up to null sets), so the orbit identity transfers.
\end{proof}

The literal-torus statements have two limitations.  No differential operator is transported, so on the literal
$\T^3$ divergence-freeness is membership of the datum in the transported subspace
$\Lp^2_\sigma$, and the uniqueness of the solution with datum $B_{\rm in}^\eps$ is stated in
the product model only (Lemma~\ref{lem:dynamics:physical-uniqueness}); and the norm statements
are stated against the trajectory's own value at time zero rather than against a separately
transported initial field, so that every norm in them is a norm of
$\Lp^2(\T^3;\R^3)$
(\texttt{FastDynamo/Analysis/LiteralThreeTorus.lean}).

\begin{theorem}[The strengthened main theorem]
  \label{thm:main-strengthened}
  \uses{thm:main-spectral, thm:main-growth-rate, def:main:witnesses,
    lem:main:normalisation, lem:main:velocity-regularity, lem:main:literal-torus,
    lem:main:literal-torus-solution,
    prop:global-spectral-transfer, lem:transfer:all-cells}
  \lean{FastDynamo.MainTheoremStrengthened,
    FastDynamo.main_theorem_strengthened,
    FastDynamo.MainTheoremStrengthened.growthRate,
    FastDynamo.MainTheoremStrengthened.growthRate_le,
    FastDynamo.MainEveryAdmissibleCellClauses,
    FastDynamo.mainEveryAdmissibleCellClauses_holds,
    FastDynamo.MainSemigroupOrbitClauses,
    FastDynamo.mainSemigroupOrbitClauses_holds}
  \leanok
  At the explicit witnesses $u$, $\eps_0$, $\gamma_0$ of
  Definition~\ref{def:main:witnesses} the following hold.
  \begin{enumerate}
    \item \emph{Explicit constants.}  $u$ is divergence-free and non-zero,
      $\eps_0=\delta_+k_*^2/K_0^2$ with $k_*=1$, and $\gamma_0>0$ satisfies both
      $\gamma_0\geq\tfrac34(448-4\pi^2)$, with $N=2\pi$, and the numeral bound
      $\gamma_0\geq288$.
    \item \emph{The sixteen physical clauses} $\mathrm{MPC}(u,\eps_0,\gamma_0)$
      of Definition~\ref{def:main:physical-clauses}.
    \item \emph{The growth-rate clause} \eqref{eq:main-growth-conclusion}, in
      the form of Theorem~\ref{thm:main-growth-rate}: for every
      $0<\eps\leq\eps_0$ the growth rate of the induction semigroup on
      $\Lp^2_\sigma(\T^3;\CC^3)$ satisfies $\gamma(u,\eps)\geq\gamma_0$, and at
      every positive diffusivity $\gamma(u,\eps)\leq\norm{\nabla_\perp u}_\infty$.
    \item \emph{Every admissible cell.}
      Proposition~\ref{prop:global-spectral-transfer} holds not only at the
      canonically chosen matching cell but at \emph{every} one: for every
      $0<\eps\leq\eps_0$ and every cell index $n$ with
      $\eps/\ell_n^2\in I$, the axial index $K_n$ is non-zero and carries an
      eigenpair, of the reduced and of the auxiliary operator, with
      $\re\lambda\geq\gamma_0$; and the Riesz-projection form of the
      proposition holds at every such cell as well.
    \item \emph{Normalisation} of the spectral clauses: the sixteen physical clauses hold
      for the named witness $\widetilde u=u/\max\{1,\norm u\}$, with
      $\norm{\widetilde u}\leq1$ and thresholds $M^{-1}\eps_0$, $M^{-1}\gamma_0$ (the whole of
      Remark~1.6 of the paper is Lemma~\ref{lem:main:normalisation}).
    \item \emph{Velocity regularity}, in the form of
      Lemma~\ref{lem:main:velocity-regularity}, together with the
      identification of the torus-side representative with the
      $\Z^2$-periodisation of the chart representative at every point of the
      plane.
    \item \emph{The literal three-torus}, in the form of
      Lemmas~\ref{lem:main:literal-torus} and~\ref{lem:main:literal-torus-solution}.
    \item \emph{The semigroup orbit of the real datum.}  For every
      $0<\eps\leq\eps_0$ the datum $B_{\rm in}^\eps$ lies in $\Lp^2_\sigma$, its orbit
      under the semigroup of Theorem~\ref{thm:dynamics:solenoidal-semigroup} is the
      explicit curve $B^\eps$ at every $t\geq0$, and that orbit obeys
      \eqref{eq:main-thm-growth}; every forward strong solution of the curl-form equation
      with datum $B_{\rm in}^\eps$ is that orbit, obeys \eqref{eq:main-thm-growth} and is
      bounded below by $\e^{t\gamma_0}\norm{B_{\rm in}^\eps}_2$; and at the same triple
      $(K,\lambda,b)$ the sixteen physical clauses of
      Definition~\ref{def:main:physical-clauses} hold.  This is the definite
      article of ``\emph{its} solution'' in Theorem~\ref{thm:main}.
  \end{enumerate}
\end{theorem}

\begin{proof}
  \leanok
  \uses{def:transfer:parameter-choice}
  Each item is the corresponding statement above, read at the named witnesses
  instead of at anonymous existentials.  Item (4) is the same closure
  certificate as in Theorem~\ref{thm:main-spectral}, except that the four
  defining properties of the parameter chain
  (Definition~\ref{def:transfer:parameter-choice}) are recorded one by one, so
  the certificate is available at every admissible diffusivity and every
  matching cell rather than only at the canonically chosen one.
\end{proof}

The structure is \texttt{FastDynamo.MainTheoremStrengthened} and the proof is
\texttt{FastDynamo.main\_theorem\_strengthened}
(\texttt{FastDynamo/MainTheoremStrengthened.lean}); the witnesses are named in
\texttt{FastDynamo/Global/MainTheoremWitnesses.lean}; the growth-rate item is
the pair of fields \texttt{growthRate} and \texttt{growthRate\_le}, and the semigroup-orbit
item is the field \texttt{semigroupOrbit}.  The
growth-rate clause, which Theorem~\ref{thm:main-spectral} does not contain, is included here,
and so is the uniqueness of the solution with datum $B_{\rm in}^\eps$ (item 8).  The dynamical
statements are Proposition~\ref{prop:main-lagrangian}; they are in a separate structure,
\texttt{MainTheoremComplete} (Theorem~\ref{thm:main-complete}), because the module that names
the velocity cannot import the flow without a cycle.

The items are separate fields, and the items stated for every $0<\eps\leq\eps_0$ choose their
triples $(K,\lambda,b)$ separately: items (2), (5), (7) and (8) each carry their own
existential over $(K,\lambda,b)$, item (5) for the normalised velocity.  Item (8) contains the sixteen physical clauses at its own triple, so the
eigenmode, the real datum, the strong solution, the semigroup orbit and the uniqueness clause
hold at one $(K,\lambda,b)$.  The triple of item (7) is the one of item (8) by the proof of
\texttt{mainLiteralTorusClauses\_holds}, not by its statement; item (2) is the existential
\texttt{MainPhysicalClauses} with a separately chosen triple; and item (4) is a statement about
every admissible cell, not about a chosen triple.
